\documentclass[11pt,a4paper]{article}
% preamble.tex — estilo común de todos los apuntes Froth.
% Cada apunte hace % preamble.tex — estilo común de todos los apuntes Froth.
% Cada apunte hace % preamble.tex — estilo común de todos los apuntes Froth.
% Cada apunte hace \input{preamble} justo después de \documentclass.
\usepackage[margin=2.4cm]{geometry}
\usepackage{xcolor}
\definecolor{litblue}{RGB}{31,79,121}
\definecolor{litgreen}{RGB}{27,94,32}
\definecolor{litorange}{RGB}{230,81,0}
\usepackage{parskip}
\usepackage{titlesec}
\titleformat{\section}{\Large\bfseries\color{litblue}}{\thesection}{0.6em}{}
\titleformat{\subsection}{\large\bfseries\color{litblue}}{\thesubsection}{0.6em}{}
\usepackage{tcolorbox}
\usepackage{fancyhdr}
\pagestyle{fancy}
\fancyhf{}
\fancyhead[L]{\small\color{litblue}Froth — Apuntes de ML}
\fancyhead[R]{\small\thepage}
\renewcommand{\headrulewidth}{0.4pt}
\usepackage[colorlinks=true,linkcolor=litblue,urlcolor=litblue]{hyperref}

% Cabecera de cada apunte: \cabecera{numero}{titulo}
\newcommand{\cabecera}[2]{%
  \begin{tcolorbox}[colback=litblue,coltext=white,colframe=litblue,arc=2mm]
    {\Large\bfseries Apunte #1 — #2}\\[3pt]
    {\small Proyecto Froth · aprender Machine Learning construyendo}
  \end{tcolorbox}\vspace{0.8em}}

% Cajas temáticas
\newtcolorbox{concepto}[1]{colback=blue!4,colframe=litblue,title={Concepto: #1},arc=1.5mm}
\newtcolorbox{practica}{colback=green!5,colframe=litgreen,title={Pruébalo tú (teclado en mano)},arc=1.5mm}
\newtcolorbox{ojo}{colback=orange!8,colframe=litorange,title={Ojo / error típico},arc=1.5mm}
\newtcolorbox{resumen}{colback=gray!8,colframe=gray!60,title={En una frase},arc=1.5mm}

\newcommand{\codigo}[1]{\texttt{#1}}
 justo después de \documentclass.
\usepackage[margin=2.4cm]{geometry}
\usepackage{xcolor}
\definecolor{litblue}{RGB}{31,79,121}
\definecolor{litgreen}{RGB}{27,94,32}
\definecolor{litorange}{RGB}{230,81,0}
\usepackage{parskip}
\usepackage{titlesec}
\titleformat{\section}{\Large\bfseries\color{litblue}}{\thesection}{0.6em}{}
\titleformat{\subsection}{\large\bfseries\color{litblue}}{\thesubsection}{0.6em}{}
\usepackage{tcolorbox}
\usepackage{fancyhdr}
\pagestyle{fancy}
\fancyhf{}
\fancyhead[L]{\small\color{litblue}Froth — Apuntes de ML}
\fancyhead[R]{\small\thepage}
\renewcommand{\headrulewidth}{0.4pt}
\usepackage[colorlinks=true,linkcolor=litblue,urlcolor=litblue]{hyperref}

% Cabecera de cada apunte: \cabecera{numero}{titulo}
\newcommand{\cabecera}[2]{%
  \begin{tcolorbox}[colback=litblue,coltext=white,colframe=litblue,arc=2mm]
    {\Large\bfseries Apunte #1 — #2}\\[3pt]
    {\small Proyecto Froth · aprender Machine Learning construyendo}
  \end{tcolorbox}\vspace{0.8em}}

% Cajas temáticas
\newtcolorbox{concepto}[1]{colback=blue!4,colframe=litblue,title={Concepto: #1},arc=1.5mm}
\newtcolorbox{practica}{colback=green!5,colframe=litgreen,title={Pruébalo tú (teclado en mano)},arc=1.5mm}
\newtcolorbox{ojo}{colback=orange!8,colframe=litorange,title={Ojo / error típico},arc=1.5mm}
\newtcolorbox{resumen}{colback=gray!8,colframe=gray!60,title={En una frase},arc=1.5mm}

\newcommand{\codigo}[1]{\texttt{#1}}
 justo después de \documentclass.
\usepackage[margin=2.4cm]{geometry}
\usepackage{xcolor}
\definecolor{litblue}{RGB}{31,79,121}
\definecolor{litgreen}{RGB}{27,94,32}
\definecolor{litorange}{RGB}{230,81,0}
\usepackage{parskip}
\usepackage{titlesec}
\titleformat{\section}{\Large\bfseries\color{litblue}}{\thesection}{0.6em}{}
\titleformat{\subsection}{\large\bfseries\color{litblue}}{\thesubsection}{0.6em}{}
\usepackage{tcolorbox}
\usepackage{fancyhdr}
\pagestyle{fancy}
\fancyhf{}
\fancyhead[L]{\small\color{litblue}Froth — Apuntes de ML}
\fancyhead[R]{\small\thepage}
\renewcommand{\headrulewidth}{0.4pt}
\usepackage[colorlinks=true,linkcolor=litblue,urlcolor=litblue]{hyperref}

% Cabecera de cada apunte: \cabecera{numero}{titulo}
\newcommand{\cabecera}[2]{%
  \begin{tcolorbox}[colback=litblue,coltext=white,colframe=litblue,arc=2mm]
    {\Large\bfseries Apunte #1 — #2}\\[3pt]
    {\small Proyecto Froth · aprender Machine Learning construyendo}
  \end{tcolorbox}\vspace{0.8em}}

% Cajas temáticas
\newtcolorbox{concepto}[1]{colback=blue!4,colframe=litblue,title={Concepto: #1},arc=1.5mm}
\newtcolorbox{practica}{colback=green!5,colframe=litgreen,title={Pruébalo tú (teclado en mano)},arc=1.5mm}
\newtcolorbox{ojo}{colback=orange!8,colframe=litorange,title={Ojo / error típico},arc=1.5mm}
\newtcolorbox{resumen}{colback=gray!8,colframe=gray!60,title={En una frase},arc=1.5mm}

\newcommand{\codigo}[1]{\texttt{#1}}

\begin{document}
\cabecera{42}{Cuando el programa no se cae sino se CUELGA: watchdogs y subprocesos}

\section{El accidente real (madrugada del 14 de julio)}

Lanzamos la re-cosecha sin techo de los 22 topics. El batch iba perfecto: 18 corpus
gigantes (4.000--16.681 papers cada uno) en unas 3 horas\ldots{} y de repente, silencio.
A la mañana siguiente: \textbf{9 horas congelado} en el topic \#20, sin error, sin
mensaje, sin avanzar. El proceso seguía ``vivo'' (77 MB de RAM, cero CPU), esperando
una respuesta de red que jamás llegó.

\section{La lección central: un cuelgue NO es un crash}

Nuestro \texttt{run\_topic} ya era ``a prueba de fallos'': envolvía todo el pipeline en
\texttt{try/except}, y cada fallo se convertía en una fila \texttt{FAILED} del reporte
para seguir con el siguiente topic. Sonaba blindado. No lo estaba.

\begin{concepto}{crash vs.\ cuelgue}
Un \textbf{crash} lanza una excepción: \texttt{try/except} la atrapa y la vida sigue.
Un \textbf{cuelgue} (hang) no lanza NADA: el programa se queda esperando para siempre
(un socket muerto, un cómputo desbocado, un driver bloqueado). Ningún
\texttt{try/except} del mundo interrumpe una espera infinita --- para eso hace falta
alguien MIRANDO DESDE AFUERA con un reloj: un \textbf{watchdog}.
\end{concepto}

Y como los topics corrían en serie, un solo cuelgue congeló el batch entero. La
probabilidad no era despreciable: 22 topics $\times$ 4 APIs $\times$ miles de requests.

\section{El arreglo: proceso hijo + reloj de pared}

Cada topic corre ahora en un \textbf{subproceso} propio
(\texttt{python -m froth.batch --one "<título>"}) vigilado por el padre con
\texttt{subprocess.run(..., timeout=TOPIC\_TIMEOUT\_S)} (2.400\,s, generoso: el topic
legítimo más lento tardó 38 min). Si el hijo se pasa del reloj, Windows lo MATA, el
padre escribe \texttt{FAILED: timed out} y sigue con el siguiente. El cuelgue pasó de
``catástrofe silenciosa de 9 horas'' a ``una fila fea en el CSV''.

De paso acotamos los dos cómputos que podían desbocarse de verdad:
el barrido DBCV para la granularidad automática es $\mathcal{O}(n^2)$ y se repite
18 veces --- sobre 16.000 puntos son horas; ahora, pasado un tope
(\texttt{DBCV\_SUBSAMPLE\_CAP} = 5.000) se barre sobre una submuestra aleatoria
(la ELECCIÓN de granularidad es estable; el costo no), y el clustering final sí usa
todos los puntos. Y el rellenado de abstracts, el único bucle de red sin presupuesto,
recibió su límite de tiempo.

\section{Bonus 1: el plan B de Beauvoir}

El topic ``Recovery of by-products of lithium from the a rare-metal granite'' cosechó
\textbf{0 papers}: sus términos derivados giraban sobre ``a rare-metal granite'', demasiado
nicho para las APIs. Arreglo: si la cosecha vuelve VACÍA, se reintenta UNA vez con una
query amplia de alto recall (las primeras palabras de contenido del título:
\texttt{recovery by-products lithium}). Resultado real: de 178 papers viejos a
\textbf{4.937}.

\section{Bonus 2: el enemigo externo (WDAC en olas)}

La política corporativa de control de aplicaciones (WDAC) bloquea a ratos las DLL
nativas de scipy/hdbscan/torch (``Una directiva de Control de aplicaciones bloqueó este
archivo''). Medido en vivo: dejó trabajar 19 minutos a un hijo y 5 minutos después
bloqueó a los otros dos en el import. Contra una política corporativa no se pelea:
se ESPERA con paciencia automatizada --- un sondeo barato
(\texttt{python -c "import hdbscan, umap, torch"}) cada 5 minutos, y cuando la ola
pasa, se lanza el trabajo pesado. El orquestador padre, además, solo importa
\texttt{config} + \texttt{pandas}: sin DLLs científicas, WDAC no puede tumbarlo.

\begin{ojo}
Tres capas distintas de defensa, porque son tres fallos distintos: \texttt{try/except}
atrapa \emph{crashes}, el watchdog con timeout mata \emph{cuelgues}, y el sondeo de
ventanas espera \emph{bloqueos del entorno}. Confundirlas es creer que un chaleco
salvavidas te protege de un rayo.
\end{ojo}

\begin{resumen}
El batch se congeló 9 horas porque un cuelgue no lanza excepción y nada lo vigilaba.
Arreglo en tres capas: subproceso por topic con timeout de pared (watchdog), cómputo
acotado (DBCV por submuestra), y sondeo de ventanas WDAC. Resultado: 22/22 topics con
estado claro y $\sim$169.000 papers cosechados; Beauvoir pasó de 178 a 4.937 con el
plan B de términos, a synthetic-ore topic de 237 a 9.723. Un sistema robusto no es el que nunca falla:
es el que convierte cualquier fallo en una fila del reporte y sigue.
\end{resumen}

\end{document}
